\section{Discussion}
\label{sec:discussion}

Taken together, the empirical results support a layered interpretation of Brazilian equity dependencies, consistent with the multi-scale perspective advocated by Raddant and Di Matteo~\cite{raddant_dimatteo_2023}. The first layer is a market-mode component. It is reflected in the moderate positive average correlation ($\bar{\rho}\approx0.24$), the increase in rolling average correlation around stress periods and the very large leading eigenvalue ($\lambda_1=21.65$, approximately 37.3\% of the trace of the correlation matrix). This component is economically interpretable as broad market synchronization. In the Brazilian context, such synchronization may reflect common exposure to domestic monetary conditions, exchange-rate dynamics, country risk and global commodity cycles, although these channels are not separately identified in the present analysis. The component is also methodologically important because it can dominate raw networks and obscure weaker local structures.

The second layer consists of candidate group and sector structure. Within-sector correlations are higher than between-sector correlations on average, and the RMT decomposition identifies four additional eigenvalues above the random upper edge. Eigenvector loadings and group-mode matrices suggest that commodity, financial, telecommunications and utility-related structures contribute to the non-random part of the spectrum. This interpretation is descriptive rather than causal: the analysis identifies structure in the dependency matrix, not the economic mechanisms that generate each edge.

The third layer is noise-compatible residual dependence. Not every correlation is economically meaningful, and not every edge retained by a graph should be interpreted as a stable economic relationship. This is why the paper combines RMT, heatmaps, dendrograms and network metrics. A result is more credible when it appears consistently across several representations rather than only in one visualization.

The network results demonstrate that market-mode removal changes topology in a structured way. Original networks have stronger average edge correlations, while group-mode networks have much lower average edge correlations. However, the group-mode PMFG has higher clustering (0.6653 versus 0.5273). This combination is a central finding: lower average edge strength does not mean absence of structure. It means that the dominant common component has been removed and localized mesoscopic organization becomes easier to observe \cite{tumminello_2005_tool,tumminello_2010_correlation}. The PMFG is particularly suited for this analysis because it preserves richer topology than the MST, retaining 168 edges with 166 triangles and 55 4-cliques that allow hub-rank shifts, clustering changes and subsector dependency maps to be studied in detail.

Clustering summarizes hierarchical organization, but it does not identify graph-theoretic hubs, bridges or local cycles. The transition from dendrograms to filtered financial networks provides complementary information. The hub reorganization after RMT filtering is particularly informative: original hubs such as GGBR4 and BBDC4 are central in networks that include the common component, while group-mode hubs such as GUAR3 occupy central positions after the leading component is removed. This suggests that centrality in financial networks should be interpreted conditionally on the matrix used to construct the network.

The forecasting results should be interpreted with appropriate caution. They do not show that machine learning uniformly dominates classical volatility models. Instead, they show that network-augmented models can approach strong baselines and that selected PMFG variables have nonzero importance in Random Forest models. The persistence of GARCH(1,1), with $\alpha+\beta$ consistently above 0.98, indicates that autoregressive volatility dynamics remain the dominant forecasting signal for Brazilian equities. Network features complement rather than replace this signal. The appropriate interpretation is incremental: topology adds structural information from the cross-sectional geometry of the market to the temporal dynamics of individual asset variance.

The broader contribution is methodological. The paper connects econophysics filtering, financial networks and volatility forecasting in a single reproducible pipeline for B3 equities. By documenting each step transparently and evaluating network features against demanding econometric benchmarks, the analysis provides a template that can be extended to other emerging markets, different time scales and more sophisticated forecasting architectures.
